% !TEX program = xelatex
% UTF-8; fill this single source from the lecture and real measurements.
\RequirePackage{iftex}
\RequireXeTeX
\documentclass[UTF8,a4paper,11pt,fontset=fandol,scheme=plain]{ctexart}
\setmainfont{Latin Modern Roman}
\setsansfont{Latin Modern Sans}
\setmonofont{Latin Modern Mono}
\usepackage[margin=25mm]{geometry}
\usepackage{amsmath}
\usepackage{graphicx}
\usepackage{eso-pic}
\usepackage[numbers,super,sort&compress]{natbib}
\linespread{1.5}
\setcounter{secnumdepth}{2}
\renewcommand{\thesection}{\arabic{section}}
\renewcommand{\thesubsection}{\thesection.\arabic{subsection}}
\ctexset{section={name={,}},subsection={name={,}}}
% natbib's public heading hook keeps real references as numbered section 5.
\renewcommand{\bibsection}{\section{参考文献}}

\newcommand{\ReportTitle}{实验题目}
\newcommand{\ReportBlank}[1]{\makebox[#1]{\hrulefill}}
\newcommand{\ReportTODO}[1]{\textbf{待补充：}#1}
\newif\ifReportLogo
\IfFileExists{assets/bnu_logo.png}{\ReportLogotrue}{\ReportLogofalse}
\ifReportLogo
  % First page only; 70mm wide, 20mm left of the text margin; no flow space.
  \AddToShipoutPictureFG*{\AtPageUpperLeft{%
    \put(\LenToUnit{5mm},\LenToUnit{-8mm}){%
      \raisebox{-\height}{\includegraphics[width=7cm]{assets/bnu_logo.png}}}}}
\else
  \PackageWarningNoLine{lab-report}{Logo missing: assets/bnu_logo.png; draft only}
\fi

\title{\ReportTitle}
\author{作者：\ReportBlank{2.4cm}\quad 学号：\ReportBlank{3.0cm}\\
        指导老师：\ReportBlank{2.0cm}\quad 实验时间：\ReportBlank{2.6cm}}
\date{}

\begin{document}
\maketitle
\noindent\ReportTODO{当前为模板草稿，请根据讲义与真实记录填写。}
\ifReportLogo\else
  \par\noindent\ReportTODO{首页校徽缺失，请补充 assets/bnu_logo.png。}
\fi

\section*{摘要}
\ReportTODO{用100—200字说明目的、方法、主要结果与结论；有真实结果时给出关键数值及单位。}
\par\noindent\textbf{关键词：}\ReportTODO{填写3—5个中文关键词，用分号分隔。}

\section{引言}\label{sec:introduction}
\ReportTODO{第一段：实验背景、研究对象及意义。}

\ReportTODO{第二段：本次实验目的、测量对象与实验思路。}

\section{原理}\label{sec:principle}
\subsection{实验原理}
\ReportTODO{必要原理、公式、物理量定义、单位及适用条件。}
\subsection{实验装置}
\ReportTODO{装置组成、连接方式及必要示意图。}
\subsection{实验方法}
\ReportTODO{实际操作、测量条件、数据处理方法。}

\section{结果与分析讨论}\label{sec:results}
\ReportTODO{由真实数据生成图表，先写结果，再分析规律、物理解释及误差；保留异常值与来源。}

\section{实验总结}\label{sec:conclusion}
\ReportTODO{主要结果、结论及至少两条有依据的改进建议。}

\section{参考文献}
\ReportTODO{填写3—6条实际查阅且可核实的文献；正文用 cite 命令作上标顺序引用。}
% Once real references are available, replace the preceding section and TODO
% with thebibliography (natbib supplies numbered section 5 via bibsection):
% \begin{thebibliography}{99}
% \bibitem{verified-source} <actual source, formatted to GB/T 7714-2015>
% \end{thebibliography}
% Do not insert invented sources or leave sample bibitems in a finished report.

\section*{Abstract}
\ReportTODO{忠实翻译中文摘要，不增加没有依据的结果。}
\par\noindent\textbf{Keywords: }\ReportTODO{与中文关键词对应的英文关键词。}

% Optional unnumbered appendix: programs, derivations or other non-data material.
% \section*{附录：程序与补充说明}
\end{document}
